% The social preview of the GitHub repository (Logo/social_preview.png, 1280 x 640 px; uploaded by hand in the settings
% of the repository): the logo, the name and the first panel of the README. Built by make_readme_figure.sh, after the
% gallery (\gallery is the PDF of readme_gallery.tex) and the logo (python3 make_logo.py ../Logo).
\documentclass[border=0pt]{standalone}
\usepackage[T1]{fontenc}
\usepackage[sc]{mathpazo}
\usepackage{tikz,graphicx}
\providecommand{\gallery}{readme_gallery.pdf}
\definecolor{accent}{HTML}{1F4E79}
\definecolor{muted}{HTML}{59636E}
\definecolor{edge}{HTML}{D0D7DE}
\begin{document}
\begin{tikzpicture}[x=1bp, y=1bp]
  \useasboundingbox (0, 0) rectangle (640, 320);
  \fill[white] (0, 0) rectangle (640, 320);
  \fill[accent] (0, 0) rectangle (640, 5);
  % left: logo, name, what it is; right: the first panel of the README, on a card; both centred vertically, the card
  % 280bp high, so that the 2:1 picture keeps a margin of about 40 px above and below it
  \node[anchor=west, inner sep=0pt] at (40, 163) {\begin{minipage}{262bp}\raggedright
    \includegraphics[height=66bp]{../Logo/FunctionalDiagrams_logo.pdf}\par\vskip 22bp
    {\color{accent}\fontsize{23}{28}\selectfont\bfseries FunctionalDiagrams}\par\vskip 12bp
    {\fontsize{14}{18}\selectfont Feynman diagrams from\\ functional derivatives}\par\vskip 9bp
    {\color{muted}\fontsize{10.5}{14}\selectfont A Mathematica package for perturbation theory,\\ effective actions and many-body physics}
  \end{minipage}};
  \node[anchor=east, inner sep=5bp, draw=edge, line width=0.6bp, rounded corners=4bp, fill=white] at (608, 163)
    {\includegraphics[height=270bp, page=1]{\gallery}};
\end{tikzpicture}
\end{document}
